针对山区洪涝救援中的地形起伏、物资不可拆分、装备有限和通信受限问题，本文建立“安全运输能力—双资源调度—连续通信保障—独立分组配置”的递进模型。采用统一投影与DEM完整像元相交确定航段高度，结合逐段载荷、交接完成时刻、能耗和两阶段充电，生成可用于四问的任务属性。

对于问题一，枚举可行箱组并采用整数规划和动态规划交叉求解。在20\%返航余量下，18架次方案达到解析下界，总能耗59.131296\,kWh、累计作业时间546.266929\,min。安全余量提高到23.08835\%前原组批仍可行；越过临界值后，最少架次数逐级增加。

对于问题二，联合决定运输任务、飞机、电池和开始时刻，得到23架次零迟到方案，最后返航5693.231489\,s、能耗66.214460\,kWh，全局有效上下界间隙4.5541\%。补充20次相同预算搜索及充电、作业延迟、能耗偏差试验，区分原日程与时刻修复：充电时间增加30\%后，顺延修复在6108.392219\,s完成；原任务能量边界仅容许约0.121877\%统一正偏差。

对于问题三，将直连失效转化为连续服务区间，联合安排中继位置、高度、提供方及资源，并在固定离散结构后精修时刻。工期优先方案含23个运输和4个中继架次，联合最后返航5836.969929\,s、总能耗68.939416\,kWh。新增完整区间检验表明，原方案对附加损耗和上线延迟敏感；0.25\,dB保障备选增加29.170597\,s工期，12份延迟修复方案均通过独立核验。

对于问题四，冻结联合日程并压缩为3个不可拆任务块，利用资源峰值、构造链与匹配核验完整枚举。两组方案仅缺1架C型运输机，三组缺7件指定型号资源。库存扰动与分量支配分析进一步证明：在固定日程和目标顺序下，两组选择不随库存变化而改变。结果说明，救援方案应在统一物理与时间约束下兼顾名义性能、扰动容忍范围和恢复代价，不能以资源总量或一次验证通过代替系统执行保障。
